% Part: first-order-logic
% Chapter: syntax-and-semantics
% Section: formation-sequences

\documentclass[../../../include/open-logic-section]{subfiles}

\begin{document}

\olfileid{fol}{syn}{fseq}

\olsection{Urutan Pambentukan}

Nemtokake !!{formula}s liwat definisi induktif, lan teknik pelengkap
kanggo mbuktekake sipat-sipat !!{formula}s liwat induksi, iku cara kang
endah lan efisien. Nanging, bisa uga migunani kanggo ndeleng
pambentukan !!{formula}s saka unsur paling dhasar, saka tahap siji menyang tahap sabanjure.
Ing kene kita nindakake iki nganggo konsep \emph{urutan pambentukan}.
%
Kanggo nuduhake carane term lan !!{formula}s bisa dikenalake kanthi
cara iki tanpa nyandhak definisi induktife, dhisik kita ngenalake
konsep rerangken simbol sembarang kang dijupuk saka sawijining
basa~$\Lang L$.

\begin{defn}[Rerangken simbol]
\ollabel{defn:string}
Upamana $\Lang L$ iku basa logika sepisanan. Sawijining
\emph{rerangken~$\Lang L$} yaiku urutan winates simbol-simbol
saka~$\Lang L$. Yen basa~$\Lang L$ wis cetha saka konteks, kita
kerep nyebut rerangken~$\Lang L$ mung minangka \emph{rerangken}.
\end{defn}

\begin{ex}
Kanggo saben basa logika sepisanan $\Lang L$, kabeh
!!{formula}s ing~$\Lang L$ iku rerangken~$\Lang L$, nanging kosok baline
ora mesthi bener. Umpamane, \[)(\Obj v_0\lif\lexists\] iku
rerangken~$\Lang L$ nanging dudu !!{formula} ing~$\Lang L$.
\end{ex}

\begin{defn}[Urutan pambentukan kanggo term]
\ollabel{defn:fseq-trm}
Urutan winates rerangken~$\Lang L$, yaiku $\tuple{t_0,\dotsc,t_n}$,
iku \emph{urutan pambentukan} kanggo term $t$ yen $t \ident t_n$
lan kanggo saben $i \leq n$, $t_i$ iku !!a{variable} utawa
!!a{constant}, utawa $\Lang L$ ngemot $k$-panggonan !!{function}~$f$
lan ana $m_1,\dotsc,m_k < i$ saengga
$t_i \ident f(t_{m_1},\dotsc,t_{m_k})$. Yen perlu mbedakake,
kita nyebut urutan pambentukan kanggo term minangka
\emph{urutan pambentukan term}. Cathetan koreksi sumber OLPL-067:
indeks argumen ing sumber diwiwiti saka nol nganti aritase,
kang marakake cacah argumene luwih siji tinimbang aritase;
ing kene indeks diwiwiti saka siji.
\end{defn}

\begin{ex}
Urutan
\[
    \tuple{\Obj c_0, \Obj v_0, \Atom{\Obj f^2_0}{\Obj c_0, \Obj v_0}, \Atom{\Obj f^1_0}{\Atom{\Obj f^2_0}{\Obj c_0, \Obj v_0}}}
\]
iku urutan pambentukan kanggo term $\Atom{\Obj f^1_0}{\Atom{\Obj
f^2_0}{\Obj c_0, \Obj v_0}}$, semono uga
\[
    \tuple{\Obj v_0, \Obj c_0, \Atom{\Obj f^2_0}{\Obj c_0, \Obj v_0}, \Atom{\Obj f^1_0}{\Atom{\Obj f^2_0}{\Obj c_0, \Obj v_0}}}.
\]
\end{ex}

\begin{defn}[Urutan pambentukan kanggo formula]
\ollabel{defn:fseq-frm}
Urutan winates rerangken~$\Lang L$, yaiku
$\tuple{!A_0,\dotsc,!A_n}$, iku \emph{urutan pambentukan}
kanggo~$!A$ yen $!A \ident !A_n$ lan kanggo saben $i \leq n$,
$!A_i$ iku !!{formula} atomik utawa ana $j,k < i$ lan
!!a{variable}~$x$ saengga salah siji saka syarat iki lumaku:
\begin{enumerate}
    \tagitem{prvNot}{$!A_i \ident \lnot !A_j$.}{}%
    \tagitem{prvAnd}{$!A_i \ident (!A_j \land !A_k)$.}{}%
    \tagitem{prvOr}{$!A_i \ident (!A_j \lor !A_k)$.}{}%
    \tagitem{prvIf}{$!A_i \ident (!A_j \lif !A_k)$.}{}%
    \tagitem{prvIff}{$!A_i \ident (!A_j \liff !A_k)$.}{}%
    \tagitem{prvAll}{$!A_i \ident \lforall[x][!A_j]$.}{}%
    \tagitem{prvEx}{$!A_i \ident \lexists[x][!A_j]$.}{}%
\end{enumerate}
Yen perlu mbedakake, kita nyebut urutan pambentukan kanggo
formula minangka \emph{urutan pambentukan formula}.
\end{defn}

\begin{ex}
\[
\tuple{
    \Atom{\Obj A^1_0}{\Obj v_0},
    \Atom{\Obj A^1_1}{\Obj c_1},
    (\Atom{\Obj A^1_1}{\Obj c_1} \land \Atom{\Obj A^1_0}{\Obj v_0}),
    \lexists[\Obj v_0][(\Atom{\Obj A^1_1}{\Obj c_1} \land \Atom{\Obj A^1_0}{\Obj v_0})]
}
\]
iku urutan pambentukan kanggo
$\lexists[\Obj v_0][(\Atom{\Obj A^1_1}{\Obj
c_1} \land \Atom{\Obj A^1_0}{\Obj v_0})]$, semono uga
\begin{multline*}
\tuple{
    \Atom{\Obj A^1_0}{\Obj v_0},
    \Atom{\Obj A^1_1}{\Obj c_1},
    (\Atom{\Obj A^1_1}{\Obj c_1} \land \Atom{\Obj A^1_0}{\Obj v_0}),
    \Atom{\Obj A^1_1}{\Obj c_1},\\
    \lforall[\Obj v_1][\Atom{\Obj A^1_0}{\Obj v_0}],
    \lexists[\Obj v_0][(\Atom{\Obj A^1_1}{\Obj c_1} \land \Atom{\Obj A^1_0}{\Obj v_0})]
}.
\end{multline*}
%
Kaya kang katon ing tuladha kapindho, urutan pambentukan bisa
ngemot unsur kang ora perlu: !!{formula}s kang diulang utawa ora
menehi sumbangan kanggo pambentukan.
\end{ex}

\begin{prop}\ollabel{prop:formed}
Saben !!{formula}~$!A$ ing~$\Frm[L]$ nduweni urutan pambentukan.
\end{prop}

\begin{proof}
Upamana $!A$ atomik. Mula urutan $\tuple{!A}$ iku
urutan pambentukan kanggo~$!A$.
%
Saiki upamana $!B$ lan~$!C$ nduweni urutan pambentukan
$\tuple{!B_0,\dotsc,!B_n}$ lan $\tuple{!C_0,\dotsc,!C_m}$
miturut urutane.
%
\begin{enumerate}
  \tagitem{prvNot}{Yen $!A \ident \lnot !B$,
      mula $\tuple{!B_0,\dotsc,!B_n,\lnot !B_n}$
      iku urutan pambentukan kanggo~$!A$.}{}
  \tagitem{prvAnd}{Yen $!A \ident (!B \land !C)$,
      mula $\tuple{!B_0,\dotsc,!B_n,!C_0,\dotsc,!C_m,(!B_n \land !C_m)}$
      iku urutan pambentukan kanggo~$!A$.}{}
  \tagitem{prvOr}{Yen $!A \ident (!B \lor !C)$,
      mula $\tuple{!B_0,\dotsc,!B_n,!C_0,\dotsc,!C_m,(!B_n \lor !C_m)}$
      iku urutan pambentukan kanggo~$!A$.}{}
  \tagitem{prvIf}{Yen $!A \ident (!B \lif !C)$,
      mula $\tuple{!B_0,\dotsc,!B_n,!C_0,\dotsc,!C_m,(!B_n \lif !C_m)}$
      iku urutan pambentukan kanggo~$!A$.}{}
  \tagitem{prvIff}{Yen $!A \ident (!B \liff !C)$,
      mula $\tuple{!B_0,\dotsc,!B_n,!C_0,\dotsc,!C_m,(!B_n \liff !C_m)}$
      iku urutan pambentukan kanggo~$!A$.}{}
  \tagitem{prvAll}{Yen $!A \ident \lforall[x][!B]$,
      mula $\tuple{!B_0,\dotsc,!B_n,\lforall[x][!B_n]}$
      iku urutan pambentukan kanggo~$!A$.}{}
  \tagitem{prvEx}{Yen $!A \ident \lexists[x][!B]$,
      mula $\tuple{!B_0,\dotsc,!B_n,\lexists[x][!B_n]}$
      iku urutan pambentukan kanggo~$!A$.}{}
\end{enumerate}
Miturut prinsip induksi tumrap !!{formula}s,
saben !!{formula} nduweni urutan pambentukan.
\end{proof}

Kita uga bisa mbuktekake kosok baline. Iki wigati amarga nuduhake
yen rong cara kita kanggo nemtokake formula iku ekuivalen: asil
loro-lorone padha. Iki uga ateges kita bisa mbuktekake teorema
ngenani formula nganggo induksi biasa marang dawa urutan pambentukan.

\begin{lem}
\ollabel{lem:fseq-init}
Upamana $\tuple{!A_0,\dotsc,!A_n}$ iku urutan pambentukan
kanggo~$!A_n$, lan $k \leq n$. Mula $\tuple{!A_0,\dotsc,!A_k}$
iku urutan pambentukan kanggo~$!A_k$.
\end{lem}

\begin{proof}
Latihan.
\end{proof}

\begin{prob}
Buktèkna \olref[fol][syn][fseq]{lem:fseq-init}.
\end{prob}

\begin{thm}
\ollabel{thm:fseq-frm-equiv}
$\Frm[L]$ iku himpunan kabeh rerangken~$\Lang L$ kang awujud~$!A$
saengga ana urutan pambentukan formula kanggo~$!A$.
\end{thm}

\begin{proof}
Upamana $F$ iku himpunan kabeh rerangken simbol ing basa~$\Lang L$
kang nduweni urutan pambentukan. Kita wis ndeleng ing
\olref[fol][syn][fseq]{prop:formed} yen $\Frm[L] \subseteq F$,
mula saiki kita mbuktekake kosok baline.

Upamana $!A$ nduweni urutan pambentukan
$\tuple{!A_0,\dotsc,!A_n}$. Kita mbuktekake yen
$!A \in \Frm[L]$ kanthi induksi kuwat marang~$n$.
Hipotesis induksi kita yaiku saben rerangken simbol kang nduweni
urutan pambentukan kanthi indeks pungkasan $m < n$ ana ing
$\Frm[L]$. Miturut definisi urutan pambentukan,
$!A \ident !A_n$ lan rerangken pungkasan kasebut atomik,
utawa ana $j,k < n$
saengga salah siji saka kasus ing ngisor iki lumaku:
\begin{enumerate}
    \tagitem{prvNot}{$!A \ident \lnot !A_j$.}{}%
    \tagitem{prvAnd}{$!A \ident (!A_j \land !A_k)$.}{}%
    \tagitem{prvOr}{$!A \ident (!A_j \lor !A_k)$.}{}%
    \tagitem{prvIf}{$!A \ident (!A_j \lif !A_k)$.}{}%
    \tagitem{prvIff}{$!A \ident (!A_j \liff !A_k)$.}{}%
    \tagitem{prvAll}{$!A \ident \lforall[x][!A_j]$.}{}%
    \tagitem{prvEx}{$!A \ident \lexists[x][!A_j]$.}{}%
\end{enumerate}
Saiki kita nalar miturut kasus. Yen $!A$ atomik, mula
$!A_n \in \Frm[L]$. Yen ora, upamana
$!A \ident (!A_j \land !A_k)$. Miturut
\olref[fol][syn][fseq]{lem:fseq-init},
$\tuple{!A_0,\dotsc,!A_j}$ lan $\tuple{!A_0,\dotsc,!A_k}$
iku urutan pambentukan kanggo $!A_j$ lan~$!A_k$ miturut urutane.
Amarga loro-lorone minangka bagean wiwitan sejati saka urutan
pambentukan kanggo~$!A$, indeks pungkasané luwih cilik tinimbang~$n$.
Mula miturut hipotesis induksi, $!A_j$ lan~$!A_k$ ana ing~$\Frm[L]$,
lan miturut definisi !!a{formula}, mangkono uga
$(!A_j \land !A_k)$. Kasus liyane lumaku kanthi nalar kang padha.
Cathetan koreksi sumber OLPL-068: ing langkah induksi iki,
ukuran kang dibandhingake yaiku indeks pungkasan urutan,
dudu cacah unsuré; sumber nyampurake loro ukuran kasebut.
Cathetan koreksi sumber OLPL-069: ing rong papan, sumber nganggo
subskrip basa kang beda saka basa ing teorema; ing kene subskrip
diselarasake karo basa ing teorema.
Cathetan koreksi sumber OLPL-070: tandha ekuivalensi ing sumber
diganti tandha identitas sintaktis supaya langkah iki nyatakake
identitas rerangken kaya dhaptar kasus sadurunge.
\end{proof}

Urutan pambentukan kanggo term nduweni sipat kang padha karo urutan
pambentukan kanggo !!{formula}s.

\begin{prop}
\ollabel{prop:fseq-trm-equiv}
$\Trm[L]$ iku himpunan kabeh rerangken~$\Lang L$ awujud $t$
saengga ana urutan pambentukan term kanggo~$t$.
\end{prop}

\begin{proof}
Latihan.
\end{proof}

\begin{prob}
Buktèkna \olref[fol][syn][fseq]{prop:fseq-trm-equiv}.
Pituduh: gunakna cara kang padha karo kang digunakake ing bukti
\olref[fol][syn][fseq]{thm:fseq-frm-equiv}.
\end{prob}

Ana rong jinis unsur kang ora perlu ing urutan pambentukan:
unsur kang diulang lan unsur kang ora magepokan karo pambentukan
formula utawa term. Loro-lorone bisa disingkirake kanthi ndeleng
urutan pambentukan minimal.

\begin{defn}[Urutan pambentukan minimal]
\ollabel{defn:minimal-fseq}
Urutan pambentukan $\tuple{!A_0, \ldots, !A_n}$ kanggo
formula~$!A$ iku \emph{urutan pambentukan minimal} kanggo~$!A$
yen kanggo saben urutan pambentukan liyane~$s$ kanggo~$!A$,
dawa~$s$ luwih gedhe utawa padha karo~$n+1$.

Semono uga, urutan pambentukan $\tuple{t_0, \ldots, t_n}$
kanggo term~$t$ iku \emph{urutan pambentukan minimal}
kanggo~$t$ yen kanggo saben urutan pambentukan liyane~$s$ kanggo~$t$,
dawa~$s$ luwih gedhe utawa padha karo~$n+1$.
\end{defn}

Elinga yen sawijining formula utawa term bisa duwe luwih saka
siji urutan pambentukan minimal, nanging kabeh urutan kasebut
ngemot rerangken kang persis padha.

\begin{prop}
\ollabel{prop:subformula-equivs}
Syarat-syarat iki ekuivalen:
\begin{enumerate}
    \item $!B$ iku sub-!!{formula} saka~$!A$.
    \item $!B$ muncul ing saben urutan pambentukan kanggo~$!A$.
    \item $!B$ muncul ing sawijining urutan pambentukan minimal
      kanggo~$!A$.
\end{enumerate}
\end{prop}

\begin{proof}
Latihan.
\end{proof}

\begin{prob}
Buktèkna \olref[fol][syn][fseq]{prop:subformula-equivs}.
\end{prob}

\begin{history}
Urutan pambentukan dikenalake dening Raymond Smullyan ing
buku teks \emph{First-Order Logic} \citep{Smullyan1968}.
Sipat-sipat tambahan urutan pambentukan ditemtokake dening
\citet{Zuckerman1973}.
\end{history}

\end{document}
